%====================================================================%
% Main-text tables                                                   %
% Generated by src/make_tables.py from results/. Do not edit by hand. %
% To change a number, re-run the analysis and then this script.      %
% Journal rules: caption above (<=15 words), legend below, no colour, %
% no thousands separators.                                           %
%====================================================================%

\begin{table}[h]
\caption{Discrimination and calibration under three validation designs}\label{tab:split}
\begin{tabular}{@{}llcc@{}}
\toprule
Dataset & Validation design & AUROC & ECE \\
\midrule
DILIrank & Stratified five-fold CV & 0.7284 $\pm$ 0.0223 & 0.1358 $\pm$ 0.0124 \\
 & Bemis--Murcko scaffold five-fold CV & 0.7306 $\pm$ 0.0183 & 0.1299 $\pm$ 0.0289 \\
 & Temporal split at 2010 & 0.6476 [0.555--0.736] & 0.4384 \\
\midrule
DICTrank & Stratified five-fold CV & 0.6992 $\pm$ 0.0234 & 0.1292 $\pm$ 0.0220 \\
 & Bemis--Murcko scaffold five-fold CV & 0.6775 $\pm$ 0.0355 & 0.1362 $\pm$ 0.0297 \\
 & Temporal split at 2010 & 0.6198 [0.553--0.688] & 0.2702 \\
\botrule
\end{tabular}
\footnotetext{All three designs use the same compounds, model and features within each dataset (796 DILIrank, 919 DICTrank). Cross-validation entries are the mean and standard deviation across five folds; the temporal entry gives the 95\% bootstrap interval for AUROC and a single value for ECE on the 156 and 267 held-out compounds. Scaffold cross-validation changes AUROC by +0.0022 in DILIrank and -0.0217 in DICTrank, that is in opposite directions, and leaves ECE within one standard deviation of stratified cross-validation in both.}
\end{table}

\begin{table}[h]
\caption{Murphy decomposition of the Brier score under random and temporal validation}\label{tab:brier}
\begin{tabular}{@{}lcccc@{}}
\toprule
 & \multicolumn{2}{c}{DILIrank} & \multicolumn{2}{c}{DICTrank} \\
\cmidrule(lr){2-3}\cmidrule(lr){4-5}
Term & Random & Temporal & Random & Temporal \\
\midrule
Uncertainty & 0.2405 & 0.2130 & 0.1982 & 0.2370 \\
Resolution (higher is better) & 0.0528 & 0.0191 & 0.0240 & 0.0125 \\
Reliability (lower is better) & 0.0242 & 0.2327 & 0.0274 & 0.0818 \\
Brier score & 0.2126 & 0.4266 & 0.2046 & 0.3062 \\
\midrule
Reliability share of degradation & \multicolumn{2}{c}{86.1\% [84.6--97.3]} & \multicolumn{2}{c}{82.5\% [75.5--124.0]} \\
Resolution share of degradation & \multicolumn{2}{c}{13.9\%} & \multicolumn{2}{c}{17.5\%} \\
Reliability-to-resolution ratio & \multicolumn{2}{c}{6.19} & \multicolumn{2}{c}{4.72} \\
\midrule
\multicolumn{5}{@{}l}{\textit{Sensitivity: random partitions pooled before binning}} \\
Reliability share of degradation & \multicolumn{2}{c}{90.4\%} & \multicolumn{2}{c}{91.8\%} \\
\botrule
\end{tabular}
\footnotetext{The random reference decomposes each of 20 same-sized partitions separately and averages the terms, so both columns rest on the same held-out size (156 and 267). The standard deviation across those partitions is 0.0066 and 0.0074 for reliability and 0.0132 and 0.0086 for resolution. \textbf{The uncertainty row comes from a separate 200-partition run} (the R1 control of Table~\ref{tab:prev}), because the 20-partition output does not retain it; the random column of the identity therefore closes only to 0.2119 and 0.2016 rather than exactly. Shares are the change in each model-dependent term divided by the sum of both changes; intervals are 2000-fold bootstrap percentiles. An upper bound above 100\% arises in resamples where resolution improves rather than deteriorates, which occurred in 0.9\% of DILIrank and 26.9\% of DICTrank resamples. Uncertainty is fixed by the held-out class balance and is not attributed. In the temporal columns, where an exact Brier score on the same predictions is available, the binned reconstruction and the exact score agree to within 0.002 (0.4271 against 0.4266 and 0.3050 against 0.3062). The pooled row gives the random reference twenty times the sample of the temporal column and is reported only to show the direction of that bias.}
\end{table}

\begin{table}[h]
\caption{Prevalence-matched random controls for expected calibration error}\label{tab:prev}
\begin{tabular}{@{}llccc@{}}
\toprule
Dataset & Control & ECE & Ratio & $p$ \\
\midrule
DILIrank & Temporal split at 2010 & 0.4384 [0.374--0.527] & --- & --- \\
 & (R1) Unconstrained random & 0.1309 $\pm$ 0.0285 & 3.35 & 0.0050 \\
 & (R2) Held-out rate matched & 0.2495 $\pm$ 0.0322 & 1.76 & 0.0050 \\
 & (R3) Both rates matched & 0.3211 $\pm$ 0.0316 & 1.37 & 0.0050 \\
\midrule
DICTrank & Temporal split at 2010 & 0.2702 [0.228--0.336] & --- & --- \\
 & (R1) Unconstrained random & 0.1299 $\pm$ 0.0237 & 2.08 & 0.0050 \\
 & (R2) Held-out rate matched & 0.1866 $\pm$ 0.0250 & 1.45 & 0.0050 \\
 & (R3) Both rates matched & 0.2226 $\pm$ 0.0233 & 1.21 & 0.0348 \\
\botrule
\end{tabular}
\footnotetext{Each control is averaged over 200 partitions. R2 matches the held-out positive rate to the temporal held-out set; R3 matches the training rate as well, at the cost of a smaller held-out set (92 and 189 compounds against 156 and 267). Ratio is temporal ECE divided by control ECE. One-sided $p$ uses the bias-corrected estimator $(1+r)/(1+N)$, where $r$ is the number of partitions with ECE at or above the temporal value; $r$ was 0 in 5 of the six comparisons and 6 (DICTrank R3) in the remaining one. The value 0.0050 is the lower bound available at 200 repeats, so 5 of these $p$ values are at that bound and the true values are smaller. All six survive Benjamini--Hochberg correction within this family, the largest $q$ being 0.0348.}
\end{table}

\begin{table}[h]
\caption{Recalibration strategies evaluated on the temporal held-out set}\label{tab:recal}
\begin{tabular}{@{}lcccc@{}}
\toprule
 & \multicolumn{2}{c}{DILIrank} & \multicolumn{2}{c}{DICTrank} \\
\cmidrule(lr){2-3}\cmidrule(lr){4-5}
Strategy & ECE & Gap closed & ECE & Gap closed \\
\midrule
(A) No recalibration & 0.4466 & --- & 0.2726 & --- \\
(B) Platt scaling, internal & 0.4147 & 8\% & 0.1679 & 42\% \\
(C) Isotonic regression, internal & 0.4085 & 10\% & 0.1739 & 40\% \\
(D) Prior-probability correction & 0.2959 & 38\% & 0.2438 & 12\% \\
(E) Oracle Platt scaling & 0.0482 & 100\% & 0.0234 & 100\% \\
\midrule
(B)--(C) across calibration fractions & \multicolumn{2}{c}{1.3\%--20.0\%} & \multicolumn{2}{c}{35.0\%--54.4\%} \\
\botrule
\end{tabular}
\footnotetext{Strategies (B) and (C) fit on the most recent 25\% of training compounds by approval year (160 and 163 compounds, spanning 1997--2009 and 1998--2009) and use no held-out information. Strategy (D) rescales the intercept of the full-training model using the held-out positive rate; it is not prospective and is reported to isolate the pure base-rate component. Table~\ref{tab:prior} shows that estimating that rate without held-out labels does not work, so (D) belongs with (E) as a ceiling. Strategy (E) is fitted on the held-out labels. Gap closed is the fraction of the distance between (A) and (E) that a strategy recovers. \textbf{The last row is the range over calibration fractions 0.15--0.40 and is the quantity to quote}; the single-fraction values above are one point in that range. (D) and (E) are monotone transformations of (A) and leave its AUROC unchanged at 0.6603 and 0.6045; (B) and (C) are applied to a model refitted without the calibration set, so their AUROC differs (0.6559 and 0.6403; 0.6131 and 0.6153), and isotonic regression can move AUROC by as much as 0.0613 at the smallest calibration fraction. Row (A) differs from the primary split because this pipeline sorts by approval year before fitting, which changes the rows drawn under subsampling.}
\end{table}

\begin{table}[h]
\caption{Split conformal coverage under random and temporal validation}\label{tab:conformal}
\begin{tabular}{@{}llcccc@{}}
\toprule
 &  & \multicolumn{2}{c}{DILIrank} & \multicolumn{2}{c}{DICTrank} \\
\cmidrule(lr){3-4}\cmidrule(lr){5-6}
Nominal & Coverage & Random & Temporal & Random & Temporal \\
\midrule
90\% & Marginal & 0.907 & 0.737 & 0.899 & 0.824 \\
 & Positive class & 0.953 & 0.958 & 0.982 & 0.994 \\
 & Negative class & 0.842 & 0.639 & 0.696 & 0.553 \\
\midrule
80\% & Marginal & 0.810 & 0.481 & 0.807 & 0.697 \\
 & Positive class & 0.890 & 0.917 & 0.939 & 0.988 \\
 & Negative class & 0.697 & 0.287 & 0.484 & 0.233 \\
\botrule
\end{tabular}
\footnotetext{The random reference is averaged over partitions with identical fitting, calibration and held-out sizes (calibration sets of 160 and 163, held-out sets of 156 and 267). Empty prediction sets occurred in 1 of the eight conditions, all on the random side: DICTrank at nominal 80\% at a rate of 0.003. Every temporal condition gave none. Split conformal guarantees marginal coverage only; the class-conditional rows are descriptive, which is why the random reference can itself fall below nominal within a class. The marginal shortfall under the temporal split is 0.163 and 0.076 at nominal 90\% and 0.319 and 0.103 at nominal 80\%, and in every case it falls on the negative class.}
\end{table}

\begin{table}[h]
\caption{Reversed temporal split as a control against label censoring}\label{tab:reverse}
\begin{tabular}{@{}llcccccccc@{}}
\toprule
Dataset & Direction and control & Train $n$ & Test $n$ & ECE & Slope & Ratio & Reliability share & Valid reps & $p$ \\
\midrule
DILIrank & Forward, temporal split & 640 & 156 & 0.4384 & 0.2700 & --- & --- & --- & --- \\
 & Reverse, temporal split & 156 & 640 & 0.4857 & 0.1331 & --- & --- & --- & --- \\
 & Reverse, (R1) unconstrained random & 156 & 640 & 0.1990 & 0.2406 & 2.44 & 95.0\% & 200/200 & 0.0050 \\
 & Reverse, (R2) held-out rate matched & 156 & 573 & 0.1911 & 0.2423 & 2.54 & 95.5\% & 200/200 & 0.0050 \\
 & Reverse, (R3) both rates matched & 156 & 583 & 0.3916 & 0.1745 & 1.24 & 96.0\% & 40/200 & 0.0244 \\
\midrule
DICTrank & Forward, temporal split & 652 & 267 & 0.2702 & 0.2223 & --- & --- & --- & --- \\
 & Reverse, temporal split & 267 & 652 & 0.2060 & 0.2384 & --- & --- & --- & --- \\
 & Reverse, (R1) unconstrained random & 267 & 652 & 0.1589 & 0.2714 & 1.30 & 86.1\% & 200/200 & 0.0100 \\
 & Reverse, (R2) held-out rate matched & 267 & 614 & 0.1443 & 0.2686 & 1.43 & 91.1\% & 200/200 & 0.0050 \\
 & Reverse, (R3) both rates matched & 267 & 614 & 0.1821 & 0.2517 & 1.13 & 76.0\% & 200/200 & 0.1294 \\
\botrule
\end{tabular}
\footnotetext{Training on compounds approved from 2010 onwards and testing on earlier ones makes the held-out labels the mature ones, so label immaturity cannot explain a collapse in this direction. Calibration still degrades: against a size-matched random control the ratio is 2.44 and 1.30 ($p$ = 0.0050 and 0.0100), and the reliability share, 95.0\% and 86.1\%, is higher than in the forward direction. \textbf{Three caveats belong with this table.} Under the strictest control (R3) the two datasets part company, DILIrank at 1.24 ($p$ = 0.0244) against DICTrank at 1.13 ($p$ = 0.1294, not significant), so the reversed evidence is weaker than the forward evidence of Table~\ref{tab:prev}. The DILIrank R3 control retains only 40 of 200 repeats, because matching its training positive rate of 0.3077 forces a subsample that usually fails the minimum-size guard; its $p$ therefore sits at the 0.0244 bound of 40 valid repeats rather than the 0.0050 bound used elsewhere, and the repeat count must be quoted with it. Absolute AUROC is not comparable with the forward direction because the reversed training set holds only 156 and 267 compounds, whereas the ECE ratio is comparable because each control uses the same training size. This is a control experiment, not a deployment scenario.}
\end{table}

\begin{table}[h]
\caption{Non-temporal splits that saturate covariate shift}\label{tab:nontemporal}
\begin{tabular}{@{}llcccccccc@{}}
\toprule
Dataset & Split & Test $n$ & Prevalence shift & Membership AUC & AUROC & ECE & Control ECE & Ratio & $p$ \\
\midrule
DILIrank & Temporal at 2010 & 156 & -0.3470 & 0.6835 & 0.6476 & 0.4384 & 0.2495 & 1.76 & 0.0050 \\
 & Molecular weight, upper half & 156 & -0.0679 & 1.0000 & 0.6889 & 0.1868 & 0.1431 & 1.30 & 0.0746 \\
 & Calculated logP, upper half & 156 & +0.0596 & 1.0000 & 0.6592 & 0.1789 & 0.1282 & 1.40 & 0.0348 \\
 & Fingerprint cluster hold-out & 156 & +0.0198 & 0.9827 & 0.7195 & 0.1570 & 0.1284 & 1.22 & 0.1741 \\
 & Random & 156 & -0.0600 & 0.4761 & 0.7068 & 0.1340 & 0.1410 & 0.95 & 0.5672 \\
\midrule
DICTrank & Temporal at 2010 & 267 & -0.1557 & 0.7917 & 0.6198 & 0.2702 & 0.1866 & 1.45 & 0.0050 \\
 & Molecular weight, upper half & 267 & -0.0712 & 0.9993 & 0.5528 & 0.2011 & 0.1529 & 1.31 & 0.0249 \\
 & Calculated logP, upper half & 267 & +0.0079 & 0.9990 & 0.5901 & 0.1726 & 0.1284 & 1.34 & 0.0299 \\
 & Fingerprint cluster hold-out & 267 & +0.0132 & 0.9880 & 0.7081 & 0.1255 & 0.1262 & 0.99 & 0.5075 \\
 & Random & 267 & +0.0449 & 0.4532 & 0.6834 & 0.1237 & 0.1190 & 1.04 & 0.4030 \\
\botrule
\end{tabular}
\footnotetext{Each split keeps the held-out size of the temporal split and carries its own random control matched to that split's held-out prevalence (200 repeats, the R2 design of Table~\ref{tab:prev}). Membership AUC is a five-fold cross-validated classifier for which half a compound belongs to, computed by the same procedure as the era classifier described in Methods. \textbf{The design did not achieve matched shift and is reported as an upper-bound contrast instead.} Molecular weight and logP are model inputs, so splitting on them drives membership AUC to 1.0000, and the fingerprint-cluster split reaches 0.9827 and 0.9880, against 0.6835 and 0.7917 for the temporal split. Even with shift saturated, no non-temporal ratio reaches the temporal one: the six shift-saturated ratios span 0.99--1.40, and including the two random splits the eight non-temporal ratios span 0.95--1.40, while the temporal split gives 1.76 and 1.45. The random split, with membership AUC near chance, gives 0.95 and 1.04, which confirms the measurement behaves. Across the ten rows the ratio tracks the absolute prevalence shift (Spearman rho = +0.515) more closely than membership AUC (rho = +0.255); \textbf{with ten rows no significance is claimed for either correlation and they are to be read as an ordering only}. The row T reproduces the R2 ratio of Table~\ref{tab:prev} exactly, which cross-checks the new script against the existing pipeline.}
\end{table}

\begin{table}[h]
\caption{Estimating the held-out prior without held-out labels}\label{tab:prior}
\begin{tabular}{@{}lllcccc@{}}
\toprule
Dataset & Prior estimator & Information used & $\hat{\pi}$ & Error & ECE & Gap closed \\
\midrule
DILIrank & (P0) Assume no shift: training rate & training labels only & 0.6547 & +0.3470 & 0.4466 & 0.0\% \\
 & (P1) Mean predicted probability & none & 0.7411 & +0.4334 & 0.4832 & -9.2\% \\
 & (P2) Expectation--maximisation & none & 0.8235 & +0.5158 & 0.5158 & -17.4\% \\
 & (P3) Black-box shift estimation & training labels only & 0.5645 & +0.2568 & 0.4244 & 5.6\% \\
 & (P4) Oracle: held-out rate & held-out labels & 0.3077 & +0.0000 & 0.2959 & 37.8\% \\
\midrule
DICTrank & (P0) Assume no shift: training rate & training labels only & 0.7699 & +0.1557 & 0.2726 & 0.0\% \\
 & (P1) Mean predicted probability & none & 0.8443 & +0.2301 & 0.2997 & -10.9\% \\
 & (P2) Expectation--maximisation & none & 0.9550 & +0.3408 & 0.3490 & -30.7\% \\
 & (P3) Black-box shift estimation & training labels only & 1.0000 & +0.3858 & 0.3858 & -45.4\% \\
 & (P4) Oracle: held-out rate & held-out labels & 0.6142 & +0.0000 & 0.2438 & 11.5\% \\
\botrule
\end{tabular}
\footnotetext{Only the prior estimate changes between rows; the correction applied to the intercept is the one of strategy (D) in Table~\ref{tab:recal}, and (P4) is that strategy. Gap closed is measured against the same uncorrected and oracle endpoints as Table~\ref{tab:recal}. \textbf{The result is a clear negative.} Every estimator that does not see held-out labels errs upward, by +0.1557 to +0.5158, and 5 of the 8 estimates make calibration worse rather than better. The best label-free result is 5.6\% for DILIrank and 0.0\% for DICTrank, against 37.8\% and 11.5\% for the oracle; for DICTrank the best available action is to do nothing. The common direction of the error is the point: all four estimators start from the model's own scores or hard predictions, and the temporal split has already pushed those upward (mean predicted probability 0.7321 and 0.8453), so a score-based estimate reads more positives while the true prevalence has fallen. Estimating a prior from a model whose calibration has collapsed is circular. In DICTrank the black-box confusion matrix becomes near-singular (determinant 0.0050) and the estimate saturates at 1.0000. Deployable recalibration is therefore limited to strategies (B) and (C) of Table~\ref{tab:recal}.}
\end{table}

\begin{table}[h]
\caption{Concept shift isolated by holding training size and prevalence fixed}\label{tab:concept}
\begin{tabular}{@{}llccccccccc@{}}
\toprule
Dataset & Arm & Train era & Train $n$ & Train rate & AUROC & ECE & Slope & Reliability & Ratio to A & $p$ \\
\midrule
DILIrank & (A) Era-matched, random 5-fold CV & post-2010 & 125 & 0.3077 & 0.7290 & 0.1443 & 0.4205 & 0.0305 & --- & --- \\
 & (A2) Era-matched, scaffold 5-fold CV & post-2010 & 125 & 0.3077 & 0.6925 & 0.1553 & 0.3533 & 0.0276 & --- & --- \\
 & (B) Forward, full pre-2010 training set & pre-2010 & 640 & 0.6547 & 0.6476 & 0.4384 & 0.2700 & 0.2327 & --- & --- \\
 & (C) Forward, training size matched to A & pre-2010 & 125 & 0.6550 & 0.6148 & 0.4643 & 0.1918 & 0.2526 & 3.22 & 0.0196 \\
 & (D) Forward, size and prevalence matched to A & pre-2010 & 125 & 0.3040 & 0.6177 & 0.2231 & 0.1953 & 0.0751 & 1.55 & 0.0392 \\
\midrule
DICTrank & (A) Era-matched, random 5-fold CV & post-2010 & 214 & 0.6142 & 0.5822 & 0.2165 & 0.1535 & 0.0629 & --- & --- \\
 & (A2) Era-matched, scaffold 5-fold CV & post-2010 & 214 & 0.6142 & 0.5636 & 0.2496 & 0.1241 & 0.0743 & --- & --- \\
 & (B) Forward, full pre-2010 training set & pre-2010 & 652 & 0.7699 & 0.6198 & 0.2702 & 0.2223 & 0.0818 & --- & --- \\
 & (C) Forward, training size matched to A & pre-2010 & 214 & 0.7746 & 0.5705 & 0.2908 & 0.1337 & 0.0939 & 1.34 & 0.0196 \\
 & (D) Forward, size and prevalence matched to A & pre-2010 & 214 & 0.6121 & 0.5771 & 0.2380 & 0.1467 & 0.0715 & 1.10 & 0.2549 \\
\botrule
\end{tabular}
\footnotetext{Every arm is scored on the same held-out compounds, the 156 and 267 approved from 2010 onwards; arms A and A2 use out-of-fold predictions so no compound is scored by a model trained on it. Arms A and D differ only in the era the training compounds come from, their size and positive rate being matched, so the gap between them is the concept-shift term. It is +0.0788 for DILIrank ($p$ = 0.0392, 1 of 50 controls matched arm A) and +0.0216 for DICTrank ($p$ = 0.2549, 12 of 50), so \textbf{concept shift is demonstrated for hepatotoxicity and not for cardiotoxicity}. Arm A2 separates Bemis--Murcko scaffolds across folds and costs only +0.0110 and +0.0331 of ECE, which rules out analogue leakage as the explanation for arm A. Arm B reproduces the temporal split of Table~\ref{tab:split} to four decimals, which fixes this experiment on the same pipeline. Along the path A to D to C to B the three increments sum exactly to +0.2941 and +0.0537: era +0.0788 and +0.0216, training prevalence +0.2412 and +0.0527, training size -0.0258 and -0.0206. \textbf{The increments depend on the path order}, which distributes the interactions, and the negative size term means that enlarging the training set from the matched size to the full pre-2010 set slightly improved calibration, so small training sets do not explain the collapse. The $p$ values compare each control against a reference distribution of 50 resamples and are not exchangeable permutation tests; read the ECE differences first.}
\end{table}

\begin{table}[h]
\caption{Calibration tracks the training prevalence gap while discrimination does not}\label{tab:sweep}
\begin{tabular}{@{}lcccccccc@{}}
\toprule
Dataset & Training rate & Gap & Train $n$ & AUROC & ECE & ECE SD & Slope & Reliability \\
\midrule
DILIrank & 0.2000 & -0.1077 & 262 & 0.6392 & 0.2153 & 0.0205 & 0.2656 & 0.0544 \\
 & 0.3000 & -0.0077 & 262 & 0.6629 & 0.1752 & 0.0284 & 0.3036 & 0.0419 \\
 & 0.3077$^{a}$ & +0.0000 & 262 & 0.6507 & 0.1827 & 0.0308 & 0.2750 & 0.0464 \\
 & 0.4000 & +0.0923 & 262 & 0.6668 & 0.1813 & 0.0349 & 0.3025 & 0.0531 \\
 & 0.5000 & +0.1923 & 262 & 0.6736 & 0.2238 & 0.0328 & 0.3114 & 0.0788 \\
 & 0.6000 & +0.2923 & 262 & 0.6758 & 0.3034 & 0.0288 & 0.3153 & 0.1288 \\
 & 0.6547$^{b}$ & +0.3470 & 262 & 0.6635 & 0.3602 & 0.0291 & 0.2910 & 0.1683 \\
 & 0.7000 & +0.3923 & 262 & 0.6589 & 0.4023 & 0.0267 & 0.2872 & 0.1976 \\
 & 0.8000 & +0.4923 & 262 & 0.6434 & 0.5079 & 0.0257 & 0.2632 & 0.2836 \\
\midrule
DICTrank & 0.2000 & -0.4142 & 178 & 0.6256 & 0.4441 & 0.0246 & 0.2303 & 0.2172 \\
 & 0.3000 & -0.3142 & 178 & 0.6316 & 0.3544 & 0.0278 & 0.2310 & 0.1531 \\
 & 0.4000 & -0.2142 & 178 & 0.6443 & 0.2681 & 0.0322 & 0.2509 & 0.0969 \\
 & 0.5000 & -0.1142 & 178 & 0.6424 & 0.2190 & 0.0298 & 0.2474 & 0.0650 \\
 & 0.6000 & -0.0142 & 178 & 0.6398 & 0.2045 & 0.0277 & 0.2370 & 0.0537 \\
 & 0.6142$^{a}$ & +0.0000 & 178 & 0.6398 & 0.2001 & 0.0310 & 0.2429 & 0.0513 \\
 & 0.7000 & +0.0858 & 178 & 0.6387 & 0.2189 & 0.0232 & 0.2436 & 0.0564 \\
 & 0.7699$^{b}$ & +0.1557 & 178 & 0.6245 & 0.2517 & 0.0238 & 0.2257 & 0.0717 \\
 & 0.8000 & +0.1858 & 178 & 0.6222 & 0.2697 & 0.0192 & 0.2215 & 0.0808 \\
\botrule
\end{tabular}
\footnotetext{Random splits only; no temporal information enters this experiment. The held-out set is fixed at the size and positive rate of the temporal test set, marked $^{a}$, and the training set is held at a fixed size while only its positive rate is varied; $^{b}$ marks the pre-2010 rate, so that row carries the gap the temporal split actually imposes. Each row averages 100 repeats. \textbf{ECE varies by a factor of 2.90 and 2.22 across the sweep while AUROC stays within 0.0366 and 0.0221}, which is the paper's central claim reproduced in a single controlled experiment. ECE is minimised at a gap of -0.0077 and +0.0000, that is at zero within one grid step, in both datasets. Regressed on the absolute gap, ECE rises by 0.0646 [0.063--0.066] and 0.0576 [0.056--0.059] per 0.10 of gap. The two sides are not symmetric: for DILIrank the slopes are 0.0401 below and 0.0681 above the held-out rate, against 0.0571 and 0.0501 for DICTrank, so overstating the training prevalence is the more damaging error for hepatotoxicity only. The calibration slope stays near 0.29 and 0.24 throughout while reliability rises 6.77-fold and 4.23-fold, which locates the damage in the intercept rather than the spread of the predictions. \textbf{This sweep should be read for shape and slope, not for shares}: its training set is held at 262 and 178 compounds, far below the 640 and 652 of the temporal split, which raises the zero-gap baseline and compresses the range. The shares belong to Table~\ref{tab:prev} and Table~\ref{tab:concept}.}
\end{table}

\begin{table}[h]
\caption{Recency weighting of the training set does not repair calibration}\label{tab:recency}
\begin{tabular}{@{}llcccccccccc@{}}
\toprule
Dataset & Training scheme & Train $n$ & Effective $n$ & Train rate & Gap & AUROC & ECE & 95\% CI & Slope & Reliability & Resolution \\
\midrule
DILIrank & Uniform weights (= temporal split) & 640 & 640.0 & 0.6547 & +0.3470 & 0.6476 & 0.4384 & [0.374--0.527] & 0.2700 & 0.2327 & 0.0191 \\
 & Half-life 30 years & 640 & 553.7 & 0.6893 & +0.3816 & 0.6501 & 0.4424 & [0.378--0.534] & 0.2597 & 0.2378 & 0.0172 \\
 & Half-life 20 years & 640 & 486.9 & 0.6995 & +0.3918 & 0.6493 & 0.4448 & [0.378--0.535] & 0.2590 & 0.2367 & 0.0094 \\
 & Half-life 15 years & 640 & 426.4 & 0.7057 & +0.3980 & 0.6399 & 0.4567 & [0.384--0.541] & 0.2433 & 0.2415 & 0.0141 \\
 & Half-life 10 years & 640 & 333.3 & 0.7096 & +0.4019 & 0.6458 & 0.4615 & [0.395--0.547] & 0.2674 & 0.2464 & 0.0151 \\
 & Half-life 7 years & 640 & 255.0 & 0.7047 & +0.3970 & 0.6508 & 0.4543 & [0.390--0.542] & 0.2970 & 0.2440 & 0.0202 \\
 & Half-life 5 years & 640 & 189.6 & 0.6914 & +0.3838 & 0.6375 & 0.4587 & [0.387--0.536] & 0.3130 & 0.2364 & 0.0240 \\
 & Half-life 3 years & 640 & 111.9 & 0.6590 & +0.3513 & 0.6586 & 0.4385 & [0.367--0.514] & 0.3998 & 0.2134 & 0.0207 \\
 & Truncated to last 30 years & 367 & 367.0 & 0.7384 & +0.4307 & 0.6173 & 0.5121 & [0.441--0.595] & 0.2214 & 0.2937 & 0.0101 \\
 & Truncated to last 20 years & 272 & 272.0 & 0.7574 & +0.4497 & 0.6512 & 0.5180 & [0.447--0.593] & 0.2832 & 0.2955 & 0.0251 \\
 & Truncated to last 15 years & 221 & 221.0 & 0.7421 & +0.4344 & 0.6698 & 0.5047 & [0.431--0.582] & 0.3058 & 0.2862 & 0.0167 \\
 & Truncated to last 10 years & 124 & 124.0 & 0.6774 & +0.3697 & 0.6466 & 0.4693 & [0.401--0.554] & 0.2748 & 0.2566 & 0.0148 \\
\midrule
DICTrank & Uniform weights (= temporal split) & 652 & 652.0 & 0.7699 & +0.1557 & 0.6198 & 0.2702 & [0.225--0.335] & 0.2223 & 0.0818 & 0.0125 \\
 & Half-life 30 years & 652 & 573.7 & 0.7837 & +0.1695 & 0.5975 & 0.2809 & [0.232--0.343] & 0.1956 & 0.0843 & 0.0066 \\
 & Half-life 20 years & 652 & 512.7 & 0.7884 & +0.1741 & 0.5712 & 0.2664 & [0.225--0.335] & 0.1573 & 0.0830 & 0.0034 \\
 & Half-life 15 years & 652 & 455.9 & 0.7918 & +0.1775 & 0.5557 & 0.2810 & [0.232--0.342] & 0.1244 & 0.0909 & 0.0098 \\
 & Half-life 10 years & 652 & 365.0 & 0.7957 & +0.1815 & 0.5343 & 0.2804 & [0.233--0.343] & 0.0807 & 0.0912 & 0.0038 \\
 & Half-life 7 years & 652 & 284.1 & 0.7979 & +0.1836 & 0.5047 & 0.2818 & [0.232--0.341] & 0.0242 & 0.0925 & 0.0056 \\
 & Half-life 5 years & 652 & 213.9 & 0.7994 & +0.1852 & 0.4792 & 0.2679 & [0.220--0.332] & -0.0395 & 0.0841 & 0.0050 \\
 & Half-life 3 years & 652 & 129.0 & 0.8053 & +0.1911 & 0.5163 & 0.2258 & [0.186--0.294] & 0.0438 & 0.0657 & 0.0026 \\
 & Truncated to last 30 years & 414 & 414.0 & 0.8043 & +0.1901 & 0.5413 & 0.3197 & [0.268--0.379] & 0.0844 & 0.1141 & 0.0084 \\
 & Truncated to last 20 years & 309 & 309.0 & 0.8058 & +0.1916 & 0.4938 & 0.3252 & [0.277--0.390] & -0.0147 & 0.1178 & 0.0069 \\
 & Truncated to last 15 years & 252 & 252.0 & 0.7897 & +0.1755 & 0.4933 & 0.3171 & [0.271--0.383] & -0.0079 & 0.1204 & 0.0101 \\
 & Truncated to last 10 years & 143 & 143.0 & 0.7692 & +0.1550 & 0.4870 & 0.2748 & [0.230--0.337] & -0.0185 & 0.0957 & 0.0034 \\
 & Truncated to last 5 years & 70 & 70.0 & 0.7857 & +0.1715 & 0.5267 & 0.2816 & [0.237--0.352] & 0.0582 & 0.0931 & 0.0096 \\
\botrule
\end{tabular}
\footnotetext{Weights are $0.5^{(2010-\mathrm{year})/\tau}$, so a smaller half-life $\tau$ concentrates the training set on recent approvals; uniform weights reproduce the temporal split of Table~\ref{tab:split}. Truncation keeps only the last $K$ years and is the simpler alternative a practitioner would try first. Effective $n$ is $(\sum w)^2/\sum w^2$ and shows what the weighting costs. \textbf{For DILIrank no half-life beats uniform weighting}: the best ECE over the whole sweep, 0.4384, is the uniform value itself, and truncation is worse at every window. \textbf{For DICTrank the apparent gain is bought by destroying resolution}: at $\tau$ = 3 years ECE falls from 0.2702 to 0.2258, but Murphy resolution falls from 0.0125 to 0.0026, roughly a fifth of its value, AUROC falls from 0.6198 to 0.5163 and the calibration slope from 0.2223 to 0.0438, on an effective sample of 129 rather than 652. Flattening the predictions lowers ECE and leaves a model that no longer separates the classes, which is the very error this paper argues against. Half-lives were chosen without looking at the held-out set, by an inner split of the training era (fit before 2000, select on 2000--2009); that selection rule did find the sweep optimum for DICTrank, so the procedure works even though what it selects is not useful. Read with Table~\ref{tab:recal}: the only intervention that recovers a large share of the lost calibration uses labelled recent compounds.}
\end{table}

\begin{table}[h]
\caption{Calibration degrades across the whole chemical space, not only unfamiliar compounds}\label{tab:applicability}
\begin{tabular}{@{}llcccccccc@{}}
\toprule
Dataset & Stratum & $n$ & Prevalence & AUROC & ECE & Control ECE & Ratio & 95\% CI & $p$ \\
\midrule
DILIrank & Nearest-neighbour Tanimoto, lowest third & 56 & 0.4107 & 0.5731 & 0.4254 & 0.2368 & 1.80 & [1.320--2.377] & 0.0323 \\
 & Nearest-neighbour Tanimoto, middle third & 48 & 0.2292 & 0.4816 & 0.5897 & 0.3429 & 1.72 & [1.393--2.078] & 0.0323 \\
 & Nearest-neighbour Tanimoto, highest third & 52 & 0.2692 & 0.8214 & 0.3764 & 0.3211 & 1.17 & [0.872--1.575] & 0.1935 \\
 & Scaffold seen in training set & 34 & --- & --- & --- & --- & --- & --- & underpowered \\
 & Scaffold unseen in training set & 122 & 0.3361 & 0.6022 & 0.4518 & 0.2622 & 1.72 & [1.480--2.096] & 0.0323 \\
\midrule
DICTrank & Nearest-neighbour Tanimoto, lowest third & 90 & 0.6111 & 0.5512 & 0.3085 & 0.2102 & 1.47 & [1.063--1.952] & 0.0323 \\
 & Nearest-neighbour Tanimoto, middle third & 88 & 0.5909 & 0.6031 & 0.3123 & 0.2173 & 1.44 & [1.048--1.896] & 0.0323 \\
 & Nearest-neighbour Tanimoto, highest third & 89 & 0.6404 & 0.7171 & 0.2724 & 0.1852 & 1.47 & [1.012--2.015] & 0.0323 \\
 & Scaffold seen in training set & 50 & 0.6000 & 0.7883 & 0.2820 & 0.2457 & 1.15 & [0.684--1.661] & 0.1613 \\
 & Scaffold unseen in training set & 217 & 0.6175 & 0.5820 & 0.2813 & 0.2051 & 1.37 & [1.102--1.726] & 0.0323 \\
\botrule
\end{tabular}
\footnotetext{Each stratum is compared against random splits matched to its own size and positive rate, because the prevalence differences documented in Table~\ref{tab:sweep} confound a raw comparison across strata; underpowered strata follow the rule used for the disjoint subsets ($n<40$ or minority class $<10$). \textbf{Of the 9 strata that could be computed, 7 have a ratio whose confidence interval lies entirely above one}, so the time split is worse than a matched random split whether the held-out compounds are chemically familiar or not. The similarity gradient, lowest third minus highest third of the ratio, is 0.63 [-0.037, 1.266] for DILIrank and -0.01 [-0.695, 0.661] for DICTrank, both intervals covering zero, while the discrimination gradient, highest third minus lowest third of AUROC, is 0.2484 [0.024--0.451] and 0.1656 [0.002--0.334], neither covering zero. \textbf{This pair of results must not be read as proof that the two measures behave differently}: contrasting a significant with a non-significant interval is not a test of their difference, the two quantities are on different scales, and the calibration interval still admits a gradient as large as 1.27. What the data support is narrower: discrimination varies with similarity, and no variation in the calibration ratio was detectable at these stratum sizes. The calibration intervals resample only the held-out compounds and hold the control mean fixed, so they are narrower than the true intervals.}
\end{table}

\begin{table}[h]
\caption{Under temporal validation the model is outscored by a constant}\label{tab:trivial}
\begin{tabular}{@{}llcccc@{}}
\toprule
Dataset & Predictor & Brier (temporal) & BSS (temporal) & Brier (random) & BSS (random) \\
\midrule
DILIrank & Trained model (XGBoost) & 0.4271 & -0.2809 & 0.2302 & +0.0868 \\
 & Constant = training prevalence & 0.3334 & +0.0000 & 0.2521 & +0.0000 \\
 & Constant = 0.5 & 0.2500 & +0.2502 & 0.2500 & +0.0084 \\
 & Constant = validation prevalence (oracle) & 0.2130 & +0.3611 & 0.2485 & +0.0143 \\
\midrule
DICTrank & Trained model (XGBoost) & 0.3050 & -0.1676 & 0.1811 & +0.0272 \\
 & Constant = training prevalence & 0.2612 & +0.0000 & 0.1862 & +0.0000 \\
 & Constant = 0.5 & 0.2500 & +0.0429 & 0.2500 & -0.3427 \\
 & Constant = validation prevalence (oracle) & 0.2370 & +0.0928 & 0.1842 & +0.0108 \\
\botrule
\end{tabular}
\footnotetext{Brier score, lower is better. The Brier Skill Score is $1-\mathrm{BS}_{\mathrm{model}}/\mathrm{BS}_{\mathrm{C0}}$, so a negative value means the model is worse than predicting the training prevalence for every compound. C0 is the only one of the three constants that can be computed at deployment time; C2 uses the validation labels and is an unreachable reference. \textbf{Under the temporal split the model is beaten by 6 of the 6 dataset-constant pairs, and under a random split of the same compounds it beats 6 of 6}, with features, model, hyperparameters and seed held fixed. AUROC on the same temporal split is 0.6476 and 0.6198, so ranking survives while the probabilities do not. Constants have no ranking and no AUROC.}
\end{table}

\begin{table}[h]
\caption{Calibration-error ratios exclude one in every control design}\label{tab:ratioci}
\begin{tabular}{@{}llcccc@{}}
\toprule
Dataset & Control design & Control ECE & Ratio & 95\% CI & ECE difference 95\% CI \\
\midrule
DILIrank & Random (standard) & 0.1309 & 3.35 & [2.878--4.005] & [0.246--0.392] \\
 & Random, validation prevalence matched & 0.2495 & 1.76 & [1.511--2.097] & [0.127--0.273] \\
 & Random, both prevalences matched & 0.3211 & 1.37 & [1.174--1.628] & [0.056--0.202] \\
\midrule
DICTrank & Random (standard) & 0.1299 & 2.08 & [1.754--2.594] & [0.098--0.207] \\
 & Random, validation prevalence matched & 0.1866 & 1.45 & [1.225--1.806] & [0.042--0.150] \\
 & Random, both prevalences matched & 0.2226 & 1.21 & [1.024--1.516] & [0.005--0.114] \\
\botrule
\end{tabular}
\footnotetext{The ratio is the temporal-split ECE divided by the mean ECE of 200 matched random-split controls. Intervals are 2.5--97.5 percentiles of 4,000 bootstrap resamples of the validation set; within each resample the control repetitions are resampled as well, so both the numerator and the denominator contribute sampling error. \textbf{All six lower bounds exceed one}, the smallest being 1.024 for DICTrank under the design that matches training and validation prevalence simultaneously. The ECE differences likewise exclude zero. Including the denominator's variability widens the intervals by a factor of at most 1.012, so the width is set almost entirely by the validation sample size (156 and 267 compounds).}
\end{table}

\begin{table}[h]
\caption{Binning-free metrics give the same ordering as expected calibration error}\label{tab:metrics}
\begin{tabular}{@{}llcccc@{}}
\toprule
Dataset & Control design & ECE & Debiased ECE & ICI & E90 \\
\midrule
DILIrank & Random (standard) & 3.35 & 6.76 & 4.23 & 3.20 \\
 & Random, validation prevalence matched & 1.76 & 2.51 & 1.88 & 1.42 \\
 & Random, both prevalences matched & 1.37 & 1.66 & 1.43 & 1.20 \\
\midrule
DICTrank & Random (standard) & 2.08 & 2.81 & 2.25 & 1.63 \\
 & Random, validation prevalence matched & 1.45 & 1.79 & 1.57 & 1.34 \\
 & Random, both prevalences matched & 1.21 & 1.35 & 1.28 & 1.09 \\
\botrule
\end{tabular}
\footnotetext{Each entry is the temporal-split value divided by the mean of 200 matched random-split controls. ECE uses ten equal-width bins. The debiased ECE subtracts the mean ECE of 1,000 null samples drawn as $y^{*}\sim\mathrm{Bernoulli}(p)$ with the predictions held fixed, which are perfectly calibrated by construction, so the subtracted quantity is the finite-sample bias of ECE. ICI and E90 are the mean and the 90th percentile of $|p-f(p)|$ from a loess calibration curve and use no bins at all. \textbf{All 24 entries exceed one}, and the debiased ratio is the larger of the two in 6 of 6 designs, because the bias is a larger share of the small control error than of the large temporal error. The bias is 14.7\% of the temporal ECE for DILIrank and 15.9\% for DICTrank, and the observed values exceed the 97.5th percentile of their null distributions in both, so the loss is not a small-sample artefact. A Spiegelhalter test rejects calibration for the controls as well ($z=4.31$ to $12.15$): the claim is a difference in degree, not a contrast between calibrated and uncalibrated.}
\end{table}

\begin{table}[h]
\caption{Tuning lowers calibration error but leaves the split-design gap}\label{tab:hparam}
\begin{tabular}{@{}llccccc@{}}
\toprule
Dataset & Configuration & Temporal ECE & Control ECE (R1) & Ratio (R1) & Control ECE (R3) & Ratio (R3) \\
\midrule
DILIrank & Fixed (as used throughout) & 0.4384 & 0.1294 & 3.39 & 0.3247 & 1.35 \\
 & Selected on training era, log loss & 0.4103 & 0.0837 & 4.90 & 0.3128 & 1.31 \\
\midrule
DICTrank & Fixed (as used throughout) & 0.2702 & 0.1272 & 2.12 & 0.2200 & 1.23 \\
 & Selected on training era, log loss & 0.2128 & 0.0888 & 2.40 & 0.1742 & 1.22 \\
\botrule
\end{tabular}
\footnotetext{A grid of 54 configurations crossing tree depth, number of trees, learning rate and minimum child weight was evaluated in full. Configurations were selected on the most recent quarter of the training era, never on the validation set. The fixed configuration used throughout this work ranks 39 of 54 for DILIrank and 39 of 54 for DICTrank by temporal ECE, so the absolute values reported elsewhere are not the best attainable and the grid spans 0.3700--0.4753 and 0.1716--0.3215. \textbf{The ratio to matched random-split controls nevertheless survives tuning} and widens against the standard control, because tuning lowers the control error more than the temporal error. Choosing with hindsight on the validation set itself, 0 and 0 of 54 configurations reach the control level, so the failure is not an artefact of how configurations were chosen. Temporal ECE rises monotonically with depth, number of trees and learning rate in both datasets.}
\end{table}

\begin{table}[h]
\caption{Minimum detectable improvement for each proposed remedy}\label{tab:power}
\begin{tabular}{@{}lllccccc@{}}
\toprule
Dataset & Remedy & Availability & $\Delta$ECE & 95\% CI & MDE & MDE / gap & Verdict \\
\midrule
DILIrank & Uncalibrated (calibration set withheld) & reference & +0.0156 & [-0.027, 0.043] & 0.0494 & 16.1\% & no useful gain \\
 & Platt, training-era calibration set & deployable & -0.0234 & [-0.069, 0.004] & 0.0533 & 17.3\% & no useful gain \\
 & Isotonic, training-era calibration set & deployable & -0.0006 & [-0.051, 0.029] & 0.0567 & 18.4\% & no useful gain \\
 & Prior-shift correction & needs test prevalence & -0.1342 & [-0.179, -0.102] & 0.0559 & 18.2\% & improves \\
 & Recency weighting at selected half-life & deployable & +0.0040 & [-0.025, 0.035] & 0.0427 & 13.9\% & no useful gain \\
 & Class weighting to validation prevalence & needs test prevalence & -0.0880 & [-0.128, -0.051] & 0.0543 & 17.7\% & improves \\
\midrule
DICTrank & Uncalibrated (calibration set withheld) & reference & +0.0233 & [-0.017, 0.045] & 0.0439 & 31.3\% & no useful gain \\
 & Platt, training-era calibration set & deployable & -0.0987 & [-0.137, -0.074] & 0.0463 & 33.0\% & improves \\
 & Isotonic, training-era calibration set & deployable & -0.0898 & [-0.138, -0.067] & 0.0511 & 36.4\% & improves \\
 & Prior-shift correction & needs test prevalence & -0.0273 & [-0.061, -0.012] & 0.0346 & 24.7\% & improves \\
 & Recency weighting at selected half-life & deployable & -0.0444 & [-0.084, 0.002] & 0.0614 & 43.8\% & undecided \\
 & Class weighting to validation prevalence & needs test prevalence & -0.0345 & [-0.059, -0.001] & 0.0405 & 28.9\% & improves \\
\botrule
\end{tabular}
\footnotetext{Each remedy is compared with the uncalibrated model on the same validation set by a paired bootstrap of 4,000 resamples in which the predictions are held fixed, so the two arms remain correlated. The gap a remedy must close is the uncalibrated ECE minus the mean ECE of matched random-split controls (0.3076 for DILIrank, 0.1403 for DICTrank). MDE is $2.80\times\mathrm{SE}(\Delta\mathrm{ECE})$, the smallest improvement that would have reached significance at two-sided $\alpha=0.05$ with 80\% power. \textbf{7 of the 12 arms were powered to detect an improvement of a quarter of the gap}, so the negative findings are not an artefact of sample size; 1 arm remains undecided and it is the one with the largest MDE. Of the 6 remedies that need no knowledge of the validation set, 2 improve calibration significantly, both in DICTrank; none does in DILIrank.}
\end{table}

\begin{table}[h]
\caption{Importance weighting by a domain classifier does not restore calibration}\label{tab:covshift}
\begin{tabular}{@{}llcccccc@{}}
\toprule
Dataset & Training weights & Effective $n$ & Max $w$ & AUROC & ECE & Resolution \\
\midrule
DILIrank & No weighting (as used throughout) & 640.0 & 1.00 & 0.6476 & 0.4384 & 0.0191 \\
 & Density ratio, uncapped & 24.6 & 72.55 & 0.6566 & 0.4755 & 0.0278 \\
 & Density ratio, capped at 20 & 26.4 & 61.27 & 0.6671 & 0.4557 & 0.0156 \\
 & Density ratio, capped at 10 & 42.0 & 36.13 & 0.6553 & 0.4689 & 0.0184 \\
 & Density ratio, capped at 5 & 61.4 & 21.25 & 0.6557 & 0.4670 & 0.0180 \\
 & Density ratio, capped at 2 & 102.8 & 11.31 & 0.6549 & 0.4617 & 0.0144 \\
\midrule
DICTrank & No weighting (as used throughout) & 652.0 & 1.00 & 0.6198 & 0.2702 & 0.0125 \\
 & Density ratio, uncapped & 13.4 & 138.47 & 0.6225 & 0.2904 & 0.0096 \\
 & Density ratio, capped at 20 & 40.9 & 35.77 & 0.6304 & 0.3022 & 0.0131 \\
 & Density ratio, capped at 10 & 61.5 & 21.41 & 0.6374 & 0.2870 & 0.0112 \\
 & Density ratio, capped at 5 & 92.7 & 13.33 & 0.6296 & 0.2912 & 0.0148 \\
 & Density ratio, capped at 2 & 147.3 & 7.19 & 0.6285 & 0.2944 & 0.0100 \\
\botrule
\end{tabular}
\footnotetext{Training compounds are weighted by $w(x)=P(\mathrm{post}|x)/(1-P(\mathrm{post}|x))$, where the probability comes from out-of-fold predictions of a classifier trained to separate the two eras by structure alone (area under the curve 0.6835 and 0.7917). The weighting uses no outcome labels from the validation era, so unlike prior-shift correction it is in principle deployable, provided the structures of the future compounds are available. Effective sample size is $(\sum w)^2/\sum w^2$. \textbf{ECE rises in all 10 of the 10 weighted arms in both datasets}, at every cap, while the effective sample size collapses by up to 97.9\%; the density ratio has a heavy tail and spends the sample. AUROC rises slightly, so the manipulation improves ranking and degrades probabilities. Applying the same weighting to the matched controls as well leaves the ratio at 1.42 and 1.27, against 1.37 and 1.21 unweighted.}
\end{table}

\begin{table}[h]
\caption{At equal shift the temporal split still loses more calibration}\label{tab:shiftmatched}
\begin{tabular}{@{}lllcccccc@{}}
\toprule
Dataset & Split axis & $\alpha^{*}$ & Shift AUC & Target & ECE (matched) & ECE (temporal) & Ratio & SE units \\
\midrule
DILIrank & Molecular weight & 0.29 & 0.6598 & 0.6704 & 0.3295 & 0.4384 & 1.33 & 7.3 \\
 & logP & 0.33 & 0.6753 & 0.6704 & 0.3845 & 0.4384 & 1.14 & 7.2 \\
 & Random direction & 0.31 & 0.6734 & 0.6704 & 0.2809 & 0.4384 & 1.56 & 11.5 \\
\midrule
DICTrank & Molecular weight & 0.46 & 0.7738 & 0.7807 & 0.2189 & 0.2702 & 1.23 & 8.2 \\
 & logP & 0.46 & 0.7672 & 0.7807 & 0.2801 & 0.2702 & 0.96 & -1.8 \\
 & Random direction & 0.46 & 0.7795 & 0.7807 & 0.2005 & 0.2702 & 1.35 & 11.5 \\
\botrule
\end{tabular}
\footnotetext{Non-temporal splits are built by blending the rank along a chemical axis with a random rank, so the amount of shift can be tuned continuously; $\alpha^{*}$ is the blend whose shift matches the temporal split. Shift is measured as the cross-validated area under the curve of a classifier separating the two sides from structure alone, the same quantity reported for the temporal split elsewhere. Every split matches the temporal target to within 0.0134. The size and prevalence of both sides are matched to the temporal split, so the axis is the only remaining difference. \textbf{In 5 of the 6 comparisons the temporal split has the larger ECE}, by 7.2 to 11.5 times the repetition-level standard error; 1 comparison is a tie and none is reversed. The matched splits hold fewer compounds than the temporal validation set, which inflates the small-sample bias of ECE on their side, so the comparison is conservative. Calibration loss under temporal validation therefore does not reduce to the amount of distribution shift.}
\end{table}

\begin{table}[h]
\caption{Annual retraining with recalibration is the one remedy that works}\label{tab:prequential}
\begin{tabular}{@{}llcccccc@{}}
\toprule
Dataset & Design & $n$ & ECE & SD & AUROC & Calibration slope & Mean prediction \\
\midrule
DILIrank & Fixed model (trained once, pre-2010) & 156 & 0.4384 & --- & 0.6476 & 0.2700 & 0.7321 \\
 & Annual retraining & 156 & 0.3751 & --- & 0.6842 & 0.3709 & 0.6677 \\
 & Annual retraining and recalibration & 156 & \textbf{0.2836} & --- & 0.7218 & 0.8705 & 0.5708 \\
 & Shuffled years & 156 & 0.1170 & 0.0216 & 0.7370 & 0.5186 & 0.5921 \\
 & Shuffled years, prevalence matched & 96 & 0.2383 & 0.0231 & 0.7299 & 0.4682 & 0.5218 \\
 & Shuffled years, recalibrated & 156 & 0.0847 & 0.0247 & 0.7167 & 0.9110 & 0.5813 \\
 & Shuffled years, prevalence matched, recalibrated & 96 & 0.2335 & 0.0233 & 0.6999 & 0.7996 & 0.5363 \\
\midrule
DICTrank & Fixed model (trained once, pre-2010) & 267 & 0.2702 & --- & 0.6198 & 0.2223 & 0.8453 \\
 & Annual retraining & 267 & 0.2583 & --- & 0.6227 & 0.2527 & 0.8157 \\
 & Annual retraining and recalibration & 267 & \textbf{0.0960} & --- & 0.5943 & 0.6973 & 0.6985 \\
 & Shuffled years & 267 & 0.1330 & 0.0221 & 0.6879 & 0.3915 & 0.7672 \\
 & Shuffled years, prevalence matched & 191 & 0.1868 & 0.0192 & 0.6916 & 0.4045 & 0.7507 \\
 & Shuffled years, recalibrated & 267 & 0.0538 & 0.0121 & 0.6750 & 0.9103 & 0.7283 \\
 & Shuffled years, prevalence matched, recalibrated & 191 & 0.1158 & 0.0196 & 0.6696 & 0.9089 & 0.7177 \\
\botrule
\end{tabular}
\footnotetext{Each validation year is predicted by a model trained only on compounds approved before it, and the predictions of all years are pooled into a single evaluation set, which is the same set of compounds used throughout this work. Recalibration uses Platt scaling fitted on the most recent quarter of each training window, the same procedure as the static recalibration analysis but on a window that moves. The shuffled-year rows are controls in which the approval years are permuted, so the window sizes are preserved and the time ordering is destroyed; they are reported with and without prevalence matching and with and without recalibration, so that every arm can be read against a control that matches it on both. Paired bootstrap differences against the fixed model are -0.0681 and -0.0200 for retraining alone and -0.1655 and -0.1759 with recalibration (DILIrank and DICTrank datasets); retraining alone is not significant in DICTrank. After retraining and recalibration the error is 1.21 times the matched control for DILIrank and 0.83 times for DICTrank, against 1.37 and 1.21 for the static design. These figures assume that a toxicity label is available in the year of approval, which is optimistic, so they bound what periodic retraining can achieve rather than estimate it.}
\end{table}
